% !TEX root = ../main.tex
\chapter{Extrapolation to very-extreme expectile levels}
\label{ch:pool-ext}

\section*{Overview}

This chapter extrapolates the pooled intermediate-level expectile
estimator of \Cref{ch:pool-int} to very-extreme levels
$\tau_n^\prime \uparrow 1$ with
$n(1-\tau_n^\prime) \to c \in [0, \infty)$. The argument is deliberately
short: the pooled intermediate result of \Cref{thm:pool-int:main} is
fed into the plug-in extrapolation lemma
\Cref{prop:bg:plug-in-clt}, which is the thesis-level formulation of
\textcite[Proposition~C.6]{DavisonPadoanStupfler2023}. The resulting
limit is the expectile analogue of the pooled Weissman extreme-quantile
limit in \Cref{thm:bg:pool-weissman}: after division by the diverging
log-rate, only the pooled Hill estimator remains visible at first
order.

\section{Motivation}
\label{sec:pool-ext:motivation}

The intermediate-level theory of \Cref{ch:pool-int} estimates
$\expectile{\tau_n}$ at levels for which the expected number of
upper-tail observations diverges. This chapter moves to the
very-extreme regime
\[
  \tau_n^\prime \uparrow 1,
  \qquad
  n(1-\tau_n^\prime)\to c\in[0,\infty),
\]
where the target lies beyond, or at the edge of, the range supported by
the observed upper order statistics. As recalled in
\Cref{sec:bg:weissman-expectile}, the iid extreme-expectile theory of
\textcite{DaouiaGirardStupfler2018} handles this regime by
extrapolating an intermediate-level expectile estimator with a
Weissman-type power factor. Direct LAWS or direct QB estimation at
$\tau_n^\prime$ is therefore not the object of the chapter: the stable
statistical anchor remains the intermediate level $\tau_n$, and the
very-extreme target is reached by extrapolation.

The pooling protocol is unchanged. Each sample still communicates only
the two primitives inherited from \textcite{DaouiaPadoanStupfler2021}:
the local Hill estimator $\hat\gamma_j(k_j)$ and the local Weissman
quantile input $\hat q_j^\star(\,\cdot\,\mid k_j)$, together with the
asymptotic covariance information needed for pooling. No raw
observations and no expectile-specific scalar, such as a local LAWS
estimate, are introduced at this stage. Consequently the
intermediate-level anchor available for extrapolation is precisely the
canonical pooled QB estimator
\[
  \hatexpectile{\tau_n}^{\mathrm{pool}}
    (\bm\omega_\gamma,\bm\omega_q)
  =
  \psifn\!\big(\hat\gamma_n(\bm\omega_\gamma)\big)
  \hatWei_n(\tau_n\mid\bm\omega_q)
\]
from \Cref{eq:pool-int:expectile-pool}.

The role of \Cref{prop:bg:plug-in-clt}, the thesis-level version of
\textcite[Proposition~C.6]{DavisonPadoanStupfler2023}, is exactly to
make this transfer modular. Once an intermediate estimator and a
tail-index estimator satisfy a joint root-rate limit, their
Weissman-extrapolated expectile has an extreme-level limit governed
only by the tail-index component. Chapter~3 supplies that joint limit
for
$\big(\hat\gamma_n(\bm\omega_\gamma),
\hatexpectile{\tau_n}^{\mathrm{pool}}(\bm\omega_\gamma,\bm\omega_q)\big)$.
The remaining sections of the present chapter therefore do not develop
a new expectile primitive; they specify the extrapolation regime,
define the pooled Weissman-type lift, and read the limiting
distribution and optimal weights from the pooled Hill component.

\section{Extrapolation regime and input limit}
\label{sec:pool-ext:regime}

The standing assumptions of \Cref{sec:pool-int:setup} remain in force.
Thus $m$ is fixed; the $m$ samples are independent and have common
distribution function $F$; $k=\sum_{j=1}^{m}k_j$ with
$k_j\to\infty$ and $k_j/n_j\to0$; the proportionality conditions
\eqref{eq:bg:pool-rates} hold; and the intermediate anchor satisfies
\[
  L_j(\tau_n)
  =
  \log\!\frac{k_j}{n_j(1-\tau_n)}
  \;\longrightarrow\;
  L_j^\star\in\R,
  \qquad j=1,\ldots,m.
\]
The moment and log-domain conventions are also inherited: throughout
we work under $\gamma\in(0,1/2)$,
$\E|\min(X,0)|^{2+\delta}<\infty$ for some $\delta>0$,
\begin{equation}
\label{eq:pool-ext:moment-rate}
  \frac{\sqrt{k}}{\Quant_X(\tau_n)}\;\longrightarrow\;0,
\end{equation}
and all logarithms of Hill, Weissman and QB quantities are understood
on the high-probability positivity event described in
\Cref{sec:pool-int:estimator}. The only strengthening relative to the
intermediate chapter is that the second-order condition is now
assumed with
\begin{equation}
\label{eq:pool-ext:rho-strict}
  \rho<0.
\end{equation}
This strict version is part of the extrapolation theory in
\textcite[Proposition~C.6]{DavisonPadoanStupfler2023}, restated here
as \Cref{prop:bg:plug-in-clt}, and is the same strict second-order
condition that appears in the pooled Weissman theorem
\Cref{thm:bg:pool-weissman}.

The new target level is a very-extreme sequence
\begin{equation}
\label{eq:pool-ext:very-extreme-level}
  \tau_n^\prime\uparrow1,
  \qquad
  n(1-\tau_n^\prime)\to c\in[0,\infty),
  \qquad
  \frac{1-\tau_n^\prime}{1-\tau_n}\to0.
\end{equation}
Write the associated Weissman log-rate as
\begin{equation}
\label{eq:pool-ext:ell-def}
  \ell_n
  \;\coloneqq\;
  \log\!\left(\frac{1-\tau_n}{1-\tau_n^\prime}\right).
\end{equation}
Then $\ell_n\to\infty$, and the extrapolation is not allowed to be so
far into the tail that the Weissman normalisation overtakes the
available root-$k$ information:
\begin{equation}
\label{eq:pool-ext:weissman-rate}
  \frac{\sqrt{k}}{\ell_n}\;\longrightarrow\;\infty.
\end{equation}
The finite-log condition $L_j(\tau_n)\to L_j^\star$ already implies
$n(1-\tau_n)\to\infty$, because
$n_j(1-\tau_n)$ has the same order as $k_j$ for every $j$. Hence
\eqref{eq:pool-ext:very-extreme-level} is exactly the pair of regimes
required by \Cref{prop:bg:plug-in-clt}: $\tau_n$ is intermediate and
$\tau_n^\prime$ is very-extreme.

The remaining high-level conditions in \Cref{prop:bg:plug-in-clt} are
also supplied by Chapter~3. With $a_n=\sqrt{k}$, the second-order
expectile--quantile expansion \eqref{eq:bg:c-gamma-rho} and
\Cref{eq:pool-int:Lambda-bullet} give
\[
  a_n A((1-\tau_n)^{-1})
  =
  \sqrt{k}\,A((1-\tau_n)^{-1})
  \;\longrightarrow\;
  \Lambda^\bullet\in\R,
\]
while \eqref{eq:pool-ext:moment-rate} gives
$a_n/\Quant_X(\tau_n)\to0$. Finally,
\Cref{rem:bg:plugin-ratio-control} verifies the deterministic
expectile-ratio condition
\[
  \frac{\sqrt{k}}{\ell_n}
  \left\{
    \log\expectile{\tau_n^\prime}
    - \log\expectile{\tau_n}
    - \gamma\ell_n
  \right\}
  \;\longrightarrow\;0.
\]
Indeed, the $A$-term in \eqref{eq:bg:c-gamma-rho} is
$O(\sqrt{k}A((1-\tau_n)^{-1})/\ell_n)=o(1)$ because
$\ell_n\to\infty$, and the first-moment contribution is
$O(\sqrt{k}/\{\Quant_X(\tau_n)\ell_n\})=o(1)$ by
\eqref{eq:pool-ext:moment-rate}.

It remains to record the exact intermediate Gaussian input that will
be fed into the plug-in extrapolation lemma. Put
\[
  \bm D^\star \coloneqq \mathrm{diag}(\bm L^\star),
  \qquad
  \bm L^\star=(L_1^\star,\ldots,L_m^\star)^\top .
\]

\begin{proposition}[Joint input for the extrapolation step]
\label{prop:pool-ext:input-clt}
Under the assumptions stated above, for any affine weight vectors
$\bm\omega_\gamma,\bm\omega_q\in\R^m$ satisfying
$\bm\omega_\gamma^\top\bm 1=\bm\omega_q^\top\bm 1=1$,
\begin{equation}
\label{eq:pool-ext:input-clt}
  \sqrt{k}
  \begin{pmatrix}
    \hat\gamma_n(\bm\omega_\gamma)-\gamma \\[2pt]
    \hatexpectile{\tau_n}^{\mathrm{pool}}
      (\bm\omega_\gamma,\bm\omega_q)/\expectile{\tau_n}-1
  \end{pmatrix}
  \;\todist\;
  \mathcal N\!\left(
  \begin{pmatrix}
    \bm\omega_\gamma^\top\Bc \\[2pt]
    \mu^\xi(\bm\omega_\gamma,\bm\omega_q)
  \end{pmatrix},
  \begin{pmatrix}
    \bm\omega_\gamma^\top\Vc\bm\omega_\gamma
      & s_{\gamma,\xi}(\bm\omega_\gamma,\bm\omega_q) \\[2pt]
    s_{\gamma,\xi}(\bm\omega_\gamma,\bm\omega_q)
      & \sigma^\xi(\bm\omega_\gamma,\bm\omega_q)^2
  \end{pmatrix}
  \right),
\end{equation}
where $\mu^\xi$ and $\sigma^\xi$ are given in
\eqref{eq:pool-int:mu-xi} and \eqref{eq:pool-int:sigma-xi}, and
\begin{equation}
\label{eq:pool-ext:gamma-xi-cross}
  s_{\gamma,\xi}(\bm\omega_\gamma,\bm\omega_q)
  =
  m(\gamma)\,\bm\omega_\gamma^\top\Vc\bm\omega_\gamma
  +
  \bm\omega_\gamma^\top\bm D^\star\Vc\bm\omega_q .
\end{equation}
\end{proposition}

\begin{proof}
\Cref{prop:pool-int:joint-clt} gives the joint Gaussian limit of
$\hat\gamma_n(\bm\omega_\gamma)$ and
$\log\hatWei_n(\tau_n\mid\bm\omega_q)$. Apply the same first-order
delta method used in the proof of \Cref{thm:pool-int:main} to the
two-dimensional map
\[
  (g,x)
  \longmapsto
  \left(
    g,\,
    \log\psifn(g)+x-\log\expectile{\tau_n}
  \right),
\]
together with the deterministic expectile--quantile bias
\eqref{eq:pool-int:b-xi-Q-def}. The second coordinate has the
asymptotic mean and variance displayed in
\eqref{eq:pool-int:mu-xi} and \eqref{eq:pool-int:sigma-xi}. Its
covariance with the first coordinate is the covariance of the
pooled-Hill component with
$m(\gamma)$ times itself plus the pooled log-Weissman component,
which is exactly \eqref{eq:pool-ext:gamma-xi-cross} by
\eqref{eq:pool-int:Sigma}. Passing from the centred log-ratio to the
relative error does not change the limit, because
$e^x-1=x+o(x)$ uniformly for $x=O_p(k^{-1/2})$.
\end{proof}

\Cref{prop:pool-ext:input-clt} is more information than the final
very-extreme theorem will visibly retain. The extrapolated limit uses
only the first component, the pooled Hill estimator, because
\Cref{prop:bg:plug-in-clt} divides by the diverging log-rate
$\ell_n$. The joint statement is nevertheless the right object to
state here: it is exactly the hypothesis of the plug-in theorem, with
\[
  \Gamma\sim
  \mathcal N\!\big(\bm\omega_\gamma^\top\Bc,\,
                   \bm\omega_\gamma^\top\Vc\bm\omega_\gamma\big),
  \qquad
  \Delta
  \text{ equal to the second component of }
  \eqref{eq:pool-ext:input-clt}.
\]

\section{The pooled extrapolated expectile estimator}
\label{sec:pool-ext:estimator}

\todo{Define the pooled extrapolated estimator by Weissman-style lifting
of the pooled intermediate-level estimator of \Cref{ch:pool-int}:
\[
  \hatexpectile{\tau_n^\prime}^{\mathrm{pool},\star}
  \;\coloneqq\;
  \left(\frac{1-\tau_n^\prime}{1-\tau_n}\right)^{-\hat\gamma_n(\bm\omega_\gamma)}
  \hatexpectile{\tau_n}^{\mathrm{pool}}(\bm\omega_\gamma,\bm\omega_q).
\]
Emphasise that this keeps the same two primitive weight vectors as
\Cref{eq:pool-int:expectile-pool}: $\bm\omega_\gamma$ acts on the
pooled Hill estimator and $\bm\omega_q$ acts on the pooled
weighted-geometric quantile input.}

\todo{Replace the old adjusted-quantile-level alternative. The
same-level equivalence in \Cref{prop:bg:expectile-quantile} is
$\expectile{\tau}/\Quant_X(\tau)\to\psifn(\gamma)$; Chapter 4 should
not introduce
$1-(1-\tau_n^\prime)/(\gamma^{-1}-1)$ as an expectile target level.
Instead record the two legitimate comparisons. First, at matched
weights $\bm\omega_\gamma=\bm\omega_q=\bm\omega$, the identity
\[
  \hatWei_n(\tau_n^\prime\mid\bm\omega)
  =
  \left(\frac{1-\tau_n^\prime}{1-\tau_n}\right)^{
    -\hat\gamma_n(\bm\omega)}
  \hatWei_n(\tau_n\mid\bm\omega)
\]
shows that
$\hatexpectile{\tau_n^\prime}^{\mathrm{pool},\star}
=\psifn(\hat\gamma_n(\bm\omega))\hatWei_n(\tau_n^\prime\mid\bm\omega)$
exactly. Second, geometric pooling of marginal extrapolated QB
expectiles is asymptotically equivalent to this matched-weight
construction by the same Taylor argument as
\Cref{prop:pool-int:A-B-equiv}. If
$\bm\omega_\gamma\ne\bm\omega_q$, plugging
$\hatWei_n(\tau_n^\prime\mid\bm\omega_q)$ into
$\psifn(\hat\gamma_n(\bm\omega_\gamma))$ changes the exponent from
$\bm\omega_\gamma$ to $\bm\omega_q$ and generally changes the
first-order limit; do not present it as an equivalent construction.}

\section{Main extrapolation theorem}
\label{sec:pool-ext:CLT}

\todo{State the main theorem as a direct application of
\Cref{prop:bg:plug-in-clt} with $a_n=\sqrt{k}$:
\[
  \frac{\sqrt{k}}{\ell_n}
  \!\left(\frac{\hatexpectile{\tau_n^\prime}^{\mathrm{pool},\star}}
                {\expectile{\tau_n^\prime}}-1\right)
  \;\todist\;
  \mathcal N\!\left(
    \bm\omega_\gamma^\top\Bc,\;
    \bm\omega_\gamma^\top\Vc\bm\omega_\gamma
  \right).
\]
Use \Cref{sec:pool-ext:regime} to verify the hypotheses of
\Cref{prop:bg:plug-in-clt}. The proof should explicitly show the
log-decomposition
\[
  \log\frac{\hatexpectile{\tau_n^\prime}^{\mathrm{pool},\star}}
           {\expectile{\tau_n^\prime}}
  =
  \ell_n\{\hat\gamma_n(\bm\omega_\gamma)-\gamma\}
  +
  \log\frac{\hatexpectile{\tau_n}^{\mathrm{pool}}}
           {\expectile{\tau_n}}
  -
  \{\log\expectile{\tau_n^\prime}
     -\log\expectile{\tau_n}
     -\gamma\ell_n\},
\]
then divide by $\ell_n/\sqrt{k}$. The second term is
$O_p(k^{-1/2})$ by \Cref{thm:pool-int:main}, hence disappears after
multiplication by $\sqrt{k}/\ell_n$, and the deterministic third term
is negligible by \Cref{rem:bg:plugin-ratio-control}.}

\section{Optimal weights for extrapolation}
\label{sec:pool-ext:weights}

\todo{Because the asymptotic variance is governed solely by
$\hat\gamma_n(\bm\omega_\gamma)$, the variance-optimal weights for
extrapolation coincide with $\bm\omega^{\mathrm{Var}}$ of
\textcite{DaouiaPadoanStupfler2021}:
$\bm\omega^{\mathrm{Var},\star} = \Vc^{-1}\bm 1/(\bm 1^\top\Vc^{-1}\bm 1)$.
In the independent common-distribution distributed case
($R_{j,\ell}\equiv 0$), these reduce
to $\omega_j \propto k_j$ (effective-sample-size weighting), which is
the same conclusion as for extreme-quantile extrapolation.
This is exactly the variance part of \Cref{prop:bg:optimal-weights} and
the distributed simplification \eqref{eq:bg:V-opt-distrib}.}

\todo{Add the AMSE statement, which is different from Chapter 3 but
identical to Paper 3's Hill/extreme-quantile criterion. If the
asymptotic bias $\bm\omega_\gamma^\top\Bc$ is retained, minimise
\[
  \mathrm{AMSE}^{\xi,\star}(\bm\omega_\gamma)
  =
  (\bm\omega_\gamma^\top\Bc)^2
  +
  \bm\omega_\gamma^\top\Vc\bm\omega_\gamma
\]
under $\bm\omega_\gamma^\top\bm 1=1$, giving the AMSE-optimal weights
of \Cref{prop:bg:optimal-weights}(ii). Explain why the
expectile-specific Chapter 3 quantities
$\Bc^\star$, $b^{\xi/\Quant}$, $m(\gamma)$ and $\bm\omega_q$ do not
enter this first-order extreme AMSE: they are in the intermediate
component $\Delta$, which \Cref{prop:bg:plug-in-clt} removes through
division by $\ell_n$.}

\todo{Clarify the role of $\bm\omega_q$. It is not identified by the
first-order extreme-level variance or AMSE, but it still affects the
finite-sample point estimator and higher-order terms through the
intermediate anchor. The natural default is the matched choice
$\bm\omega_q=\bm\omega_\gamma$, especially when
$\bm\omega_\gamma$ is chosen as $\bm\omega^{\mathrm{Var}}$ or
$\bm\omega^{\mathrm{AMSE}}$, because matched weights recover the
Paper-3 geometric Weissman form exactly at $\tau_n^\prime$.}

\section{Bias-reduction}
\label{sec:pool-ext:bias-reduction}

\todo{Bias-reduced versions of the pooled tail index lift to
bias-reduced versions of the pooled extrapolated expectile. Following
the bias-reduction scheme of
\textcite[\S2.2]{DaouiaPadoanStupfler2021}, define
$\overline\gamma_n(\bm\omega) = \hat\gamma_n(\bm\omega) - k^{-1/2}
\bm\omega^\top \hatBc$, and propagate this into
$\hatexpectile{\tau_n^\prime}^{\mathrm{pool},\star}$:
\[
  \overline{\xi}_{\tau_n^\prime}^{\,\mathrm{pool},\star}
  =
  \left(\frac{1-\tau_n^\prime}{1-\tau_n}\right)^{
    -\overline\gamma_n(\bm\omega_\gamma)}
  \hatexpectile{\tau_n}^{\mathrm{pool}}
    (\bm\omega_\gamma,\bm\omega_q).
\]
If $\hatBc\toprob\Bc$, the centred limit has mean zero and variance
$\bm\omega_\gamma^\top\Vc\bm\omega_\gamma$. Note explicitly that no
separate correction for $\Bc^\star$ or $b^{\xi/\Quant}$ is needed at
this first-order extreme scale: those are intermediate-anchor bias
terms and vanish after the $\sqrt{k}/\ell_n$ normalisation.}

\section{Asymptotic confidence intervals}
\label{sec:pool-ext:CI}

\todo{State log-scale confidence intervals for
$\expectile{\tau_n^\prime}$ in the spirit of
\textcite[Corollary~4]{DaouiaPadoanStupfler2021}, but make the bias
calibration explicit. With estimated weights
$\hat{\bm\omega}_\gamma\toprob\bm\omega_\gamma$ satisfying exact
affine normalisation and $\hatVc\toprob\Vc$, the variance-only interval
\[
  \Prob\!\left(\expectile{\tau_n^\prime}\in
    \hatexpectile{\tau_n^\prime}^{\mathrm{pool},\star}\exp\!
    \left[\pm z_{1-\alpha/2}\log\!\Big(\frac{1-\tau_n}{1-\tau_n^\prime}\Big)
    \sqrt{\frac{\hat{\bm\omega}_\gamma^\top\hatVc\hat{\bm\omega}_\gamma}{k}}
    \right]\right)
  \to 1-\alpha.
\]
is valid under undersmoothing, for instance
\Cref{eq:pool-int:undersmoothing}, or after replacing
$\hat\gamma_n$ by the bias-reduced estimator of
\Cref{sec:pool-ext:bias-reduction}.}

\todo{Add the biased-regime version. If the asymptotic mean
$\mu^{\xi,\star}(\bm\omega_\gamma)=\bm\omega_\gamma^\top\Bc$ is
retained and consistently estimated by
$\hat\mu^{\xi,\star}=\hat{\bm\omega}_\gamma^\top\hatBc$, the log-scale
interval should be centred at
\[
  \log\hatexpectile{\tau_n^\prime}^{\mathrm{pool},\star}
  -
  \hat\mu^{\xi,\star}\frac{\ell_n}{\sqrt{k}}
  \;\pm\;
  z_{1-\alpha/2}\,
  \ell_n
  \sqrt{\frac{\hat{\bm\omega}_\gamma^\top\hatVc\hat{\bm\omega}_\gamma}{k}}.
\]
Keep the practical discussion qualitative and tied to the theorem:
the interval width scales with the extrapolation distance $\ell_n$ and
with the pooled-Hill variance. Do not introduce simulations or
empirical performance claims in this thesis.}
